\documentclass[11pt]{article}

\usepackage[margin=1in]{geometry}
\usepackage{amsmath}
\usepackage{amssymb}
\usepackage{hyperref}

\title{Stage 3D Physics Note: Nontrivial Hidden Metric and Ricci Ingredients}
\author{GL-with-AI Project}
\date{\today}

\newcommand{\GRSP}{\texttt{GR\_SPACEDIM}}
\newcommand{\CHSP}{\texttt{CH\_SPACEDIM}}

\begin{document}
\maketitle

\section{Scope}

This note summarizes the Stage 3D symbolic derivation scaffold. It is a
derivation artifact only. It does not implement CCZ4 source terms, add
\texttt{hww}/\texttt{Aww} variables, or make the current
\texttt{BlackStringToy} scaffold physically meaningful.

\section{Metric Ansatz}

The first nontrivial hidden-metric check uses the diagonal SO(3)-symmetric
spatial metric
\begin{equation}
  dl^2
  =
  a(x,z)\,dx^2
  + c(x,z)\,dz^2
  + x^2 q(x,z)\,d\theta^2
  + x^2 q(x,z)\sin^2\theta\,d\phi^2 .
\end{equation}
Here \(q(x,z)=\gamma_{ww}\) is the physical hidden cartoon component. The
factors \(x^2\) and \(x^2\sin^2\theta\) are coordinate reconstruction factors,
not part of the hidden Cartesian-like component itself.

\section{Christoffel Checks}

The key hidden-sector Christoffels are
\begin{align}
  \Gamma^x_{\theta\theta}
  &=
  -\frac{1}{2a}\partial_x(x^2q),&
  \Gamma^z_{\theta\theta}
  &=
  -\frac{1}{2c}\partial_z(x^2q),\\
  \Gamma^\theta_{x\theta}
  &=
  \frac{1}{2}\partial_x\log(x^2q)
  =
  \frac{1}{x}+\frac{1}{2}\partial_x\log q,&
  \Gamma^\theta_{z\theta}
  &=
  \frac{1}{2}\partial_z\log q .
\end{align}
The \(\phi\) versions satisfy
\begin{align}
  \Gamma^x_{\phi\phi}
  &=
  \Gamma^x_{\theta\theta}\sin^2\theta,&
  \Gamma^z_{\phi\phi}
  &=
  \Gamma^z_{\theta\theta}\sin^2\theta,\\
  \Gamma^\phi_{x\phi}
  &=
  \Gamma^\theta_{x\theta},&
  \Gamma^\phi_{z\phi}
  &=
  \Gamma^\theta_{z\theta}.
\end{align}

\section{Flat Limit and Multiplicity}

Setting \(a=c=q=1\) recovers the Stage 3C flat spherical Christoffels. For any
angular tensor satisfying
\begin{equation}
  T_{\phi\phi}=T_{\theta\theta}\sin^2\theta,
\end{equation}
the angular contractions agree:
\begin{equation}
  g^{\phi\phi}T_{\phi\phi}
  =
  g^{\theta\theta}T_{\theta\theta}.
\end{equation}
For the black-string reduction,
\begin{equation}
  N_{\rm hidden}=\GRSP-\CHSP=4-2=2,
\end{equation}
so these two equivalent angular sectors give two hidden contributions.

\section{Ricci Structural Checks}

The Stage 3D script checks the structural Ricci identities
\begin{equation}
  \frac{R_{\phi\phi}}{\sin^2\theta} - R_{\theta\theta}=0
\end{equation}
and
\begin{equation}
  g^{\theta\theta}R_{\theta\theta}
  +
  g^{\phi\phi}R_{\phi\phi}
  =
  2g^{\theta\theta}R_{\theta\theta}.
\end{equation}
It also verifies that \(R_{xx}\), \(R_{zz}\), \(R_{\theta\theta}\), and
\(R_{\phi\phi}\) vanish in the flat limit.

\section{Constant-\texorpdfstring{\(q\)}{q} Ricci Regression}

Because the general abstract Ricci expressions are large, the script also
asserts a nontrivial constant-\(q\) regression. For
\begin{equation}
  dl^2 = dx^2 + dz^2 + q_0 x^2 d\Omega_2^2,
\end{equation}
the metric is flat only when \(q_0=1\). The asserted Ricci components are
\begin{align}
  R_{xx} &= 0,&
  R_{zz} &= 0,&
  R_{xz} &= 0,\\
  R_{\theta\theta} &= 1-q_0,&
  R_{\phi\phi} &= (1-q_0)\sin^2\theta,
\end{align}
and the Ricci scalar is
\begin{equation}
  R=\frac{2(1-q_0)}{q_0 x^2}.
\end{equation}
This is a self-guarding Ricci check, not a completed CCZ4 RHS derivation.

The constant-\(q_0\) regression is a useful non-flat check, but it does not
test the derivative sector of the general \(q(x,z)\) expressions. Since
\(\partial_x q=\partial_z q=0\) and all second derivatives vanish for
\(q=q_0\), the terms involving gradients and Hessians of \(\gamma_{ww}\)
remain unguarded by an internal assertion. A later derivation stage should
introduce a non-constant profile \(q(x,z)\) with independently known curvature,
or an equivalent reference calculation, to regression-test these terms.

\section{Remaining Work}

The off-diagonal reduced component \(g_{xz}\), non-constant \(q(x,z)\)
derivative-term regression, conformal variables \(\chi,h_{ww}\), full CCZ4 RHS
terms, regularized small-\(x\) limits, C++ translation, and unit-test fixtures
remain future work.

\end{document}
